% GENERATED FILE. Edit canonical lessons, then run npm run generate:books.
% canonical-source-hash: fnv1a64:0802d4df9cd83399
% canonical-lessons: GU-C217-hava, GU-C217-phukavu, GU-C217-adhi, GU-C217-vantol, GU-C217-phelavu

\chapter{Wind and dust}
\label{ch:gu-wind-and-dust}
\hlchaptermodality{217}

\begin{chapteropening}
\textbf{By the end of this chapter:} \emph{I can say that the wind will blow and that a dust storm will come.}

Complete the last lesson of chapter 217: I can say that the wind will blow and that a dust storm will come.
\end{chapteropening}

\section[\texorpdfstring{\gu{હવા} (havā)}{હવા (havā)}]{\gu{હવા} (havā) — air}

\label{lesson:GU-C217-hava}



\begin{quote}
Before the new one: say the Gujarati for temperature, then the Gujarati for degree.
\end{quote}

\subsection*{You'll want to know: \gu{હવા}}
\textbf{\gu{હવા}} — \emph{havā} — \textquotedblleft{}air\textquotedblright{}.

\begin{quote}
\textbf{\gu{સવારે} \gu{હવા} \gu{ઠંડી} \gu{હશે}.} — \emph{savāre havā ṭhanḍī haśe} — \textquotedblleft{}The air will be cool in the morning.\textquotedblright{}
\end{quote}

\begin{cousinweb}
\textbf{\gu{હવા}} (\emph{havā}) is a loan from Arabic \emph{hawā'}, \textquotedblleft{}air\textquotedblright{}, which reached Gujarati through Persian. Its Arabic root is \emph{h-w-y} (Lane).
\end{cousinweb}

\begin{grammarlens}[title={What you've built}]
One more word for everyday life at A2.
\end{grammarlens}

\subsection*{Your turn}
\begin{itemize}
\raggedright
  \item \emph{Say it:} \emph{havā}
  \item \emph{Say it:} \emph{havā}, once more
  \item \emph{Say it:} \emph{savāre havā ṭhanḍī haśe} — The air will be cool in the morning.
  \item \emph{Recall:} say the Gujarati for temperature, then the Gujarati for degree, then say \emph{havā} again
\end{itemize}

\begin{tcolorbox}[breakable,colback=teal!4,colframe=teal!35!black,title={Before you move on}]
What does \gu{હવા} mean? (\textquotedblleft{}Air\textquotedblright{}.) Say it once more.
\end{tcolorbox}
\section[\texorpdfstring{\gu{ફૂંકાવું} (phū̃kāvũ)}{ફૂંકાવું (phū̃kāvũ)}]{\gu{ફૂંકાવું} (phū̃kāvũ) — to blow}

\label{lesson:GU-C217-phukavu}



\begin{quote}
Before the new one: say the Gujarati for degree, then the Gujarati for air.
\end{quote}

\subsection*{You'll want to know: \gu{ફૂંકાવું}}
\textbf{\gu{ફૂંકાવું}} — \emph{phū̃kāvũ} — \textquotedblleft{}to blow\textquotedblright{}.

\begin{quote}
\textbf{\gu{રાત્રે} \gu{ઠંડો} \gu{પવન} \gu{ફૂંકાશે}.} — \emph{rātre ṭhanḍo pavan phū̃kāśe} — \textquotedblleft{}A cold wind will blow at night.\textquotedblright{}
\end{quote}

\begin{cousinweb}
\textbf{\gu{ફૂંકાવું}} (\emph{phū̃kāvũ}) is built on \emph{phū̃kvũ}, \textquotedblleft{}to blow\textquotedblright{}, with the \emph{-ā-} that makes a verb happen by itself, as wind does. CDIAL lists \emph{phū̃kvũ} under \emph{*phūtka}, \textquotedblleft{}blowing\textquotedblright{}, a reconstructed form (CDIAL 9102).
\end{cousinweb}

\begin{grammarlens}[title={What you've built}]
One more word for everyday life at A2.
\end{grammarlens}

\subsection*{Your turn}
\begin{itemize}
\raggedright
  \item \emph{Say it:} \emph{phū̃kāvũ}
  \item \emph{Say it:} \emph{phū̃kāvũ}, once more
  \item \emph{Say it:} \emph{rātre ṭhanḍo pavan phū̃kāśe} — A cold wind will blow at night.
  \item \emph{Recall:} say the Gujarati for degree, then the Gujarati for air, then say \emph{phū̃kāvũ} again
\end{itemize}

\begin{tcolorbox}[breakable,colback=teal!4,colframe=teal!35!black,title={Before you move on}]
What does \gu{ફૂંકાવું} mean? (\textquotedblleft{}To blow\textquotedblright{}.) Say it once more.
\end{tcolorbox}
\section[\texorpdfstring{\gu{આંધી} (ā̃dhī)}{આંધી (ā̃dhī)}]{\gu{આંધી} (ā̃dhī) — dust storm}

\label{lesson:GU-C217-adhi}



\begin{quote}
Before the new one: say the Gujarati for air, then the Gujarati for to blow.
\end{quote}

\subsection*{You'll want to know: \gu{આંધી}}
\textbf{\gu{આંધી}} — \emph{ā̃dhī} — \textquotedblleft{}dust storm\textquotedblright{}.

\begin{quote}
\textbf{\gu{બપોરે} \gu{રણમાં} \gu{આંધી} \gu{આવશે}.} — \emph{bapore raṇmā̃ ā̃dhī āvśe} — \textquotedblleft{}A dust storm will come in the desert at midday.\textquotedblright{}
\end{quote}

A feminine noun.

\begin{cousinweb}
\textbf{\gu{આંધી}} (\emph{ā̃dhī}) is a tadbhava. CDIAL lists it, glossed \textquotedblleft{}duststorm\textquotedblright{}, under Sanskrit \emph{andhikā}, \textquotedblleft{}night\textquotedblright{}, built on \emph{andha}, \textquotedblleft{}blind\textquotedblright{}, so the storm is one that brings darkness (CDIAL 390).
\end{cousinweb}

\begin{grammarlens}[title={What you've built}]
One more word for everyday life at A2.
\end{grammarlens}

\subsection*{Your turn}
\begin{itemize}
\raggedright
  \item \emph{Say it:} \emph{ā̃dhī}
  \item \emph{Say it:} \emph{ā̃dhī}, once more
  \item \emph{Say it:} \emph{bapore raṇmā̃ ā̃dhī āvśe} — A dust storm will come in the desert at midday.
  \item \emph{Recall:} say the Gujarati for air, then the Gujarati for to blow, then say \emph{ā̃dhī} again
\end{itemize}

\begin{tcolorbox}[breakable,colback=teal!4,colframe=teal!35!black,title={Before you move on}]
What does \gu{આંધી} mean? (\textquotedblleft{}Dust storm\textquotedblright{}.) Say it once more.
\end{tcolorbox}
\section[\texorpdfstring{\gu{વંટોળ} (vanṭoḷ)}{વંટોળ (vanṭoḷ)}]{\gu{વંટોળ} (vanṭoḷ) — whirlwind}

\label{lesson:GU-C217-vantol}



\begin{quote}
Before the new one: say the Gujarati for to blow, then the Gujarati for dust storm.
\end{quote}

\subsection*{You'll want to know: \gu{વંટોળ}}
\textbf{\gu{વંટોળ}} — \emph{vanṭoḷ} — \textquotedblleft{}whirlwind\textquotedblright{}.

\begin{quote}
\textbf{\gu{ઉનાળામાં} \gu{વંટોળ} \gu{આવશે}.} — \emph{unāḷāmā̃ vanṭoḷ āvśe} — \textquotedblleft{}Whirlwinds will come in summer.\textquotedblright{}
\end{quote}

A masculine noun.

\begin{cousinweb}
The origin of \textbf{\gu{વંટોળ}} (\emph{vanṭoḷ}) is not settled. CDIAL has Sanskrit \emph{vātālī}, \textquotedblleft{}whirlwind\textquotedblright{}, built on \emph{vāta}, \textquotedblleft{}wind\textquotedblright{}, but lists no Gujarati form under it (CDIAL 11500).
\end{cousinweb}

\begin{grammarlens}[title={What you've built}]
One more word for everyday life at A2.
\end{grammarlens}

\subsection*{Your turn}
\begin{itemize}
\raggedright
  \item \emph{Say it:} \emph{vanṭoḷ}
  \item \emph{Say it:} \emph{vanṭoḷ}, once more
  \item \emph{Say it:} \emph{unāḷāmā̃ vanṭoḷ āvśe} — Whirlwinds will come in summer.
  \item \emph{Recall:} say the Gujarati for to blow, then the Gujarati for dust storm, then say \emph{vanṭoḷ} again
\end{itemize}

\begin{tcolorbox}[breakable,colback=teal!4,colframe=teal!35!black,title={Before you move on}]
What does \gu{વંટોળ} mean? (\textquotedblleft{}Whirlwind\textquotedblright{}.) Say it once more.
\end{tcolorbox}
\section[\texorpdfstring{\gu{ફેલાવું} (phelāvũ)}{ફેલાવું (phelāvũ)}]{\gu{ફેલાવું} (phelāvũ) — to spread}

\label{lesson:GU-C217-phelavu}



\begin{quote}
Before the new one: say the Gujarati for dust storm, then the Gujarati for whirlwind.
\end{quote}

\subsection*{You'll want to know: \gu{ફેલાવું}}
\textbf{\gu{ફેલાવું}} — \emph{phelāvũ} — \textquotedblleft{}to spread\textquotedblright{}.

\begin{quote}
\textbf{\gu{પવનથી} \gu{ધૂળ} \gu{ગામમાં} \gu{ફેલાશે}.} — \emph{pavanthī dhūḷ gāmmā̃ phelāśe} — \textquotedblleft{}With the wind, dust will spread through the village.\textquotedblright{}
\end{quote}

\begin{cousinweb}
\textbf{\gu{ફેલાવું}} (\emph{phelāvũ}) is a tadbhava. CDIAL lists Gujarati \emph{phelvũ} and \emph{phelāvvũ}, beside Hindi \emph{phailnā}, \textquotedblleft{}to spread\textquotedblright{}, under Sanskrit \emph{prathita}, \textquotedblleft{}spread out\textquotedblright{}, extended with \emph{-ll-} (CDIAL 8651).
\end{cousinweb}

\begin{grammarlens}[title={What you've built}]
One more word for everyday life at A2.
\end{grammarlens}

\subsection*{Your turn}
\begin{itemize}
\raggedright
  \item \emph{Say it:} \emph{phelāvũ}
  \item \emph{Say it:} \emph{phelāvũ}, once more
  \item \emph{Say it:} \emph{pavanthī dhūḷ gāmmā̃ phelāśe} — With the wind, dust will spread through the village.
  \item \emph{Recall:} say the Gujarati for dust storm, then the Gujarati for whirlwind, then say \emph{phelāvũ} again
\end{itemize}

\begin{tcolorbox}[breakable,colback=teal!4,colframe=teal!35!black,title={Before you move on}]
What does \gu{ફેલાવું} mean? (\textquotedblleft{}To spread\textquotedblright{}.) Say it once more.
\end{tcolorbox}
